% Independent analytical audit; input fragment.
\section{Regular Weyl forcing from the subsidiary fields}
\label{sec:weyl-source-energy}

This section is conditional on propagation of the stated strong subsidiary
weights.  It establishes an autonomous, regular, forced spin-two system for
physical radiation.  It does not establish the remaining geometric/gauge
reconstruction or the all-angular elliptic boundary estimate.

\subsection{Physical Ricci source and its conformal order}

Use the $C_{\rm Z4c}=0$ tensor equations with the stated Appendix-C
$\Theta$ modification.  Let $n$ be the physical unit normal, $\bar n=n/\Omega$
the conformal unit normal, and retain the physical spatial covector $Z_i=z_i$.
The linearized physical Ricci projections are
\begin{align}
 R_{nn}&=H+2\Omega^2\partial_i z^i
 -3\kappa_1(1+\kappa_2)T,\label{eq:weyl-ricci-nn}\\
 R_{ni}&=-M_i,\label{eq:weyl-ricci-ni}\\
 R_{ij}&=-2D^{\rm phys}_{(i}z_{j)}+
 \gamma^{\rm phys}_{ij}
 \left[\kappa_1(1+\kappa_2)T
 -\frac23\Omega G^kz_k\right].
 \label{eq:weyl-ricci-spatial}
\end{align}
The normal-normal expression follows from the physical Raychaudhuri equation
and the $q+2T$ evolution.  The trace-free spatial expression follows from the
$A_{ij}$ equation.  Its trace is fixed by the off-shell Gauss identity
$R_{nn}+\gamma_{\rm phys}^{ij}R_{ij}=H$.
In particular
$D^i_{\rm phys}z_i=\Omega^2\partial_i z^i-\Omega G^iz_i$;
replacing this by the conformal divergence would miss the displayed isotropic
term.  The unmodified tensor $\Theta$ equation adds $-6w_*T$ to
$R_{nn}$ and the compensating opposite spatial trace; this has the same
conformal order but is a different source.

Assume
\[
 H=\Omega^4h_4,\quad T=\Omega^4\tau_4,\quad
 M_i=\Omega^3m_i,\quad z_i=\Omega^3\zeta_i.
\]
Then $R_{ab}^{\rm phys}=\Omega^2r_{ab}$ in a regular conformal frame.
For $\kappa=\kappa_1(1+\kappa_2)$ the explicit frame components are
\begin{align*}
 r_{\bar n\bar n}&=h_4+6G\cdot\zeta+
 2\Omega\partial_i\zeta^i-3\kappa\tau_4,\\
 r_{\bar n i}&=-m_i,\\
 r_{ij}&=-2\Omega\partial_{(i}\zeta_{j)}
 -5(G_i\zeta_j+G_j\zeta_i)
 +\delta_{ij}\left[\frac43G\cdot\zeta+\kappa\tau_4\right].
\end{align*}
These expressions contain no inverse conformal factor.  To see the spatial
formula directly,
\[
 D^{\rm phys}_iz_j=\Omega^2\left[
 \Omega\partial_i\zeta_j+4G_i\zeta_j+G_j\zeta_i
 -\delta_{ij}G\cdot\zeta\right].
\]
Coordinate components are regular too, because
\[
 R_{ti}=-\frac{A}{\Omega}M_i+B^jR_{ji},\qquad
 R_{tt}=\frac{A^2}{\Omega^2}R_{nn}
 -\frac{2A}{\Omega}B^iM_i+B^iB^jR_{ij}.
\]

\subsection{Exact conformal Cotton source}

Set
\[
 \ell_{ab}=\frac12\left(r_{ab}
 -\frac16\bar g_{ab}\operatorname{tr}_{\bar g}r\right).
\]
The physical Schouten tensor is $L^{\rm phys}_{ab}=\Omega^2\ell_{ab}$,
because $R_{\rm phys}=\Omega^4\operatorname{tr}_{\bar g}r$.
Define the physical Cotton tensor without an antisymmetrization convention
ambiguity by
\[
 A^{\rm phys}_{bcd}=\nabla^{\rm phys}_c L^{\rm phys}_{db}
 -\nabla^{\rm phys}_dL^{\rm phys}_{cb}.
\]
For $d^a{}_{bcd}=\Omega^{-1}\bar C^a{}_{bcd}$ the divergence law is
\begin{equation}
 \bar\nabla_a d^a{}_{bcd}=J_{bcd},\qquad
 J_{bcd}=\Omega^{-1}A^{\rm phys}_{bcd}.
 \label{eq:weyl-divergence-source}
\end{equation}
The exact linear source is
\begin{align}
 J_{bcd}={}&\Omega(\bar\nabla_c\ell_{db}
 -\bar\nabla_d\ell_{cb})
 +3(G_c\ell_{db}-G_d\ell_{cb})\notag\\
 &-\bar g_{cb}G^e\ell_{de}
 +\bar g_{db}G^e\ell_{ce}.
 \label{eq:weyl-regular-cotton}
\end{align}
For verification, the physical/conformal connection difference gives
\[
 \nabla^{\rm phys}_c(\Omega^2\ell_{db})
 =\Omega^2\bar\nabla_c\ell_{db}
 +\Omega\big[4G_c\ell_{db}+G_d\ell_{cb}+G_b\ell_{dc}
 -\bar g_{cd}G^e\ell_{eb}-\bar g_{cb}G^e\ell_{de}\big].
\]
Antisymmetrizing in $c,d$ proves~\eqref{eq:weyl-regular-cotton}.
Every term is smooth at scri under the stated subsidiary regularity.
The one conformal derivative of $r$ requires correspondingly commuted
subsidiary estimates; the base energy alone is not claimed to provide it.

\subsection{First-order electric/magnetic evolution and null flux}

Decompose the Weyl-symmetry field $d$ into symmetric trace-free tensors
$E_{ij}=d_{i\bar n j\bar n}$ and its magnetic dual $B_{ij}$.
Choose the orientation of the magnetic dual so the flat equations read
$\partial_tE-\operatorname{curl}B=\mathcal J_E$ and
$\partial_tB+\operatorname{curl}E=\mathcal J_B$.
Here
\[
 (\operatorname{curl}Q)_{ij}
 =\epsilon_{kl(i}\partial^kQ_{j)}{}^l,
 \qquad
 (a\mathbin\times Q)_{ij}=\epsilon_{kl(i}a^kQ_{j)}{}^l.
\]
The exact sources can be specified invariantly as the STF projections of
$J_{i j\bar n}$ and of the right-dual source
$J^\star_{bcd}=\epsilon_{cd}{}^{ef}J_{bef}/2$, respectively, with the same
orientation.  This definition fixes their signs independently of coordinate
epsilon conventions.

On the CMC conformal Minkowski background, the Cartesian spatial frame is
Fermi-propagated along $\bar n$.  Its acceleration is
$a_i=\partial_i\log A=-G_i/A$ and
$\bar K_{ij}=-\delta_{ij}/(aA)$.  Consequently the evolution can be written
\begin{align}
 (\partial_t-B^k\partial_k)E
 &=A\operatorname{curl}B-\frac3aE
 +2A(a\mathbin\times B)+A\mathcal J_E,\\
 (\partial_t-B^k\partial_k)B
 &=-A\operatorname{curl}E-\frac3aB
 -2A(a\mathbin\times E)+A\mathcal J_B.
\end{align}
Equivalently these are the STF evolution projections of
\eqref{eq:weyl-divergence-source}; the latter remains the convention-independent
exact definition.  All zeroth-order frame coefficients are bounded on the
closed ball, since $A>0$ there.

The identity on the ten electric/magnetic components is a common
symmetrizer, using the Euclidean STF tensor inner product.  In direction $n$,
the symmetric curl block has eigenvalue magnitudes $1$, $1/2$, and $0$;
the vector/scalar components are retained, not projected away.  Hence its
normal operator $\mathcal M_n$ has operator norm one.  At scri,
$B^n=-A$, and therefore
\[
 B^n I+A\mathcal M_n=A(\mathcal M_n-I)\leq0.
\]
The boundary flux of the full electric/magnetic energy is nonpositive.
There is no incoming Weyl mode and no physical radiation boundary value is
freely prescribed.  The bounded connection terms and regular source give
the usual finite-time forced energy estimate.

\subsection{Divergence constraints and source integrability}

The remaining projections of~\eqref{eq:weyl-divergence-source} are the
electric and magnetic divergence constraints.  In the isotropic CMC
background their leading expressions are
\[
 \mathcal C_{E,i}=\partial^jE_{ij}-J_{\bar n i\bar n},\qquad
 \mathcal C_{B,i}=\partial^jB_{ij}-J^\star_{\bar n i\bar n},
\]
or, without frame convention dependence, the corresponding normal
projections of $\bar\nabla_a d^a{}_{bcd}-J_{bcd}$.
Their propagation has principal part
\[
 (\partial_t-B\cdot\nabla)\mathcal C_E
 =\frac A2\operatorname{curl}\mathcal C_B+\text{bounded lower terms},
\]
with the opposite curl sign for $\mathcal C_B$.
The source must satisfy its Cotton identities.  At linear order about the
flat physical background, the contracted Ricci Bianchi identity implies
$\nabla_{\rm phys}^bA^{\rm phys}_{bcd}=0$; commuting derivatives introduces
no background physical curvature.  The cyclic and trace identities follow
from the Schouten definition and the same contracted Bianchi identity.
Their conformal versions are precisely the compatibility conditions needed
for homogeneous divergence-constraint propagation.  They are consequences
of the closed subsidiary evolution and compatible initial Ricci data, not
conditions satisfied by arbitrary independently prescribed $r_{ab}$.

The divergence modes have principal speeds of magnitude at most $A/2$
relative to the normal flow, so they too are outflow at scri.  Initial Weyl
Gauss/Codazzi compatibility is still necessary.  An independently evolved
regular Bianchi field does not by itself prove that a metric reconstructed
from it solves the chosen Z4c/gauge system.

Finally, $\delta d$ is gauge invariant here because the background Weyl
tensor vanishes.  The source depends only on the propagated subsidiary
fields.  This yields the physical triangular estimate
``subsidiary source $\longrightarrow$ rescaled Weyl radiation'', independent
of lapse/shift perturbations at linear order.  The remaining geometric and
gauge reconstruction must be supplied separately.

\subsection{Exact York identity and the off-constraint electric correction}

Here set the metric-derived Dirac field $\partial^jh_{ij}$ to zero;
nonzero prescribed Dirac memory can first be separated by its own elliptic
York potential.  Put $d_0=(aS)^{-1}$, avoiding confusion with the Weyl field.
Define
\begin{align*}
 k_c&=\frac a2\Omega^2\Delta c-
 \frac a2\Omega G\cdot\nabla c-\frac{3c}{2a},&
 k_h&=\frac{3a}{2}h_{ij}G_iG_j,\\
 (a_c)_{ij}&=\frac{a\Omega}{2}(\partial_i\partial_jc)^{\rm TF}
 +\frac a2(G_i\partial_jc+G_j\partial_ic)^{\rm TF},\\
 (a_h)_{ij}&=-a C_h^k{}_{ij}G_k+ad_0h_{ij}
 -\frac{a\Omega}{2}\Delta h_{ij},
 \qquad
 C_h^k{}_{ij}=\tfrac12(\partial_i h_{jk}+\partial_jh_{ik}-\partial_kh_{ij}).
\end{align*}
The trace and momentum identities are
\[
 k=k_c+k_h-\frac a4H,\qquad
 \Omega\partial^j(a_c+a_h)_{ij}
 -2(a_c+a_h)_{ij}G^j=\frac23\partial_i(k_c+k_h).
\]
Consequently the residual $\mathcal R=a-a_c-a_h$ obeys
\[
 \partial^j(\Omega^{-2}\mathcal R_{ij})
 =\Omega^{-3}\left(M_i-\frac a6\partial_iH\right).
\]
A flat York decomposition $\mathcal R=\Omega^2(L_{\rm conf}W+T_{\rm TT})$
therefore gives
\[
 \Delta_LW=m-\frac a2G h_4-\frac a6\nabla(\Omega h_4),\qquad
 \Delta_L=\Delta+\tfrac13\nabla\operatorname{div}.
\]
The gradient term introduces no necessary derivative loss: using
$\Delta_L\nabla\phi=(4/3)\nabla\Delta\phi$, it has the particular
potential $-(a/8)\nabla\Delta^{-1}(\Omega h_4)$, with the appropriate
boundary lift.  Elliptic boundary conditions and kernel control still need
to be specified.

The free $T_{\rm TT}$ is physical information not bounded by $H,M$ alone.
Its relation to the Weyl field has an important off-constraint correction.
Let
\[
 Q_{ij}=[R^{(3)}_{ij}+K K_{ij}-K_{ik}K^k{}_j]^{\rm TF}.
\]
The Gauss decomposition gives
$E^{\rm phys}_{ij}=Q_{ij}-(1/2)(R^{\rm phys}_{ij})^{\rm TF}$.
Direct conformal spatial geometry gives
$\mathcal R=-a\Omega Q$.  Therefore, using
$(R^{\rm phys}_{ij})^{\rm TF}=-2(D^{\rm phys}_{(i}z_{j)})^{\rm TF}$,
\[
 \boxed{\quad
 \frac{\mathcal R}{\Omega^2}
 =-a E+\frac a\Omega
 (D^{\rm phys}_{(i}z_{j)})^{\rm TF} .\quad}
\]
Here $E=E^{\rm phys}/\Omega$ is the regular electric component of the
rescaled Weyl tensor used above.  The last term is $O(\Omega)$ under
$z=\Omega^3\zeta$.  Thus Weyl energy controls the missing weighted TT
curvature, with a known regular subsidiary correction.  Omitting that
correction would silently identify Gauss electric curvature with Weyl
curvature off the Einstein constraint surface.
